% Section 5: Partisan Fairness — Summary of C.5 (~1,000 words)

\subsection{Partisan Fairness as a Validation Criterion}

Critics of algorithmic redistricting sometimes argue that a neutral
algorithm cannot prevent partisan imbalance because political geography
is itself partisan. Urban concentration of Democratic voters, combined
with compactness optimization, systematically packs Democrats into
landslide districts while Republicans win many moderate districts---
a pattern documented by \citet{chen2013unintentional} and
\citet{rodden2019cities} under the name ``unintentional gerrymandering.''

This critique is empirically correct: any compact redistricting method
will produce some partisan imbalance. The relevant validation question
is not whether the algorithm produces perfect partisan symmetry (an
impossible standard) but whether it produces \emph{less} partisan
imbalance than enacted plans, and whether the remaining imbalance
reflects geography rather than deliberate manipulation. Paper C.5
\citep{dellalibera2026c5} answers both questions affirmatively.

\subsection{Efficiency Gap: Definition and Measurement}

The efficiency gap \citep{stephanopoulos2015partisan} measures
asymmetric vote waste between parties. A vote is ``wasted'' if it
is cast for a losing candidate or cast for a winning candidate beyond
the margin of victory. For each party, the efficiency gap is:
\begin{equation}
  EG = \frac{W_D - W_R}{V}
\end{equation}
where $W_D$ and $W_R$ are total wasted Democratic and Republican votes
respectively, and $V$ is total votes cast. Positive values indicate
Republican advantage; negative values indicate Democratic advantage.
The absolute value $|EG|$ measures the magnitude of asymmetry
regardless of direction.

C.5 computes efficiency gaps for all 15 competitive states with at
least 3 congressional districts, across three election cycles
(2016, 2018, 2020). Algorithmic plans are compared to enacted plans
from the same states.

\subsection{Key Finding: Mean $|EG| = 0.04$ for Algorithmic Plans}

Table~\ref{tab:eg-summary} presents the main efficiency gap results
from C.5. Algorithmic plans achieve a mean efficiency gap of $-3.2\%$
(Democratic advantage, reflecting urban concentration) with mean
absolute efficiency gap $|EG| = 0.040$. Enacted plans show a mean
efficiency gap of $+5.1\%$ (Republican advantage) with mean absolute
efficiency gap $|EG| = 0.078$---nearly double.

\begin{table}[htb]
\centering
\caption{Efficiency gap summary: algorithmic vs.\ enacted plans
  (C.5, 15 competitive states, 2016--2020 elections). Mean $|EG|$
  for algorithmic plans is less than half that of enacted plans.}
\label{tab:eg-summary}
\begin{tabular}{lcccc}
\toprule
Plan type & Mean EG & SD EG & Mean $|EG|$ & States $|EG| > 7\%$ \\
\midrule
Algorithmic & $-3.2\%$ & $0.8\%$ & $0.040$ & 0 of 15 \\
Enacted     & $+5.1\%$ & $2.3\%$ & $0.078$ & 8 of 15 \\
\midrule
Difference  & $8.3\%$ & --- & $0.038$ & --- \\
\bottomrule
\end{tabular}
\end{table}

Eight of 15 enacted plans exceed the $7\%$ threshold often cited as
evidence of substantial partisan bias in state constitutional cases
\citep{lwv2018}. No algorithmic plan exceeds this threshold.

\subsection{The Compactness--Efficiency Gap Separation}

A key finding of C.5 is that enacted plans do not achieve their
partisan advantage by accepting lower compactness. In Arizona and Nevada,
algorithmic and enacted plans have \emph{identical} Polsby-Popper
scores---yet enacted plans show 6.6 and 8.4 percentage-point higher
efficiency gaps respectively. This proves that the partisan bias in
enacted plans cannot be explained as a consequence of compactness
trade-offs: mapmakers introduce bias orthogonally to geometric
optimization.

This finding has direct legal implications. The defense ``our maps
are reasonably compact, therefore any partisan bias is geographic''
is empirically refuted for these states. Courts can use algorithmic
plans as a compactness-controlled baseline: if an enacted plan achieves
similar $PP$ to the algorithmic plan but exhibits substantially higher
$|EG|$, the excess reflects manipulation, not geography.

\subsection{Multiple Fairness Metrics Converge}

C.5 examines five partisan fairness metrics (efficiency gap,
mean-median difference, partisan bias at 50\% vote share, declination,
and seats-votes elasticity). All five metrics converge on consistent
conclusions:

\begin{itemize}
  \item Algorithmic plans show modest Democratic advantage consistent
    with urban concentration ($-2$ to $-3$ pp across metrics).
  \item Enacted plans show substantial Republican advantage ($+4$ to
    $+6$ pp) reflecting deliberate cracking and packing.
  \item Cross-metric correlations between EG, partisan bias at 50\%,
    and declination exceed $r = 0.89$, confirming that findings are
    not artifacts of metric selection.
\end{itemize}

The convergence of five independent metrics on the same substantive
conclusion---roughly $8$ pp partisan gap between algorithmic and
enacted plans---substantially strengthens the validation claim.

\subsection{Temporal Consistency of Partisan Fairness}

Efficiency gaps in algorithmic plans vary by less than $0.1$ percentage
points across the 2016--2020 election cycle (coefficients of variation
below $0.05$ for all states). This temporal stability confirms that
the near-zero partisan bias of algorithmic plans is a structural
feature of geographic optimization, not a coincidence of one election
year's turnout.

Enacted plan efficiency gaps are similarly stable over time within a
redistricting cycle, but at much higher absolute levels. Pennsylvania's
enacted EG ranges from $+7.4\%$ to $+7.6\%$ across three elections;
the algorithmic range is $-2.7\%$ to $-2.9\%$. The gap is durable,
large, and stable.

\subsection{Partisan Fairness and the Impossibility Defense}

The algorithm's near-zero $|EG|$ is a \emph{byproduct} of geometric
neutrality, not a targeted optimization. METIS never accesses partisan
data; it cannot directly minimize efficiency gap. That algorithmic
redistricting produces $|EG| = 0.04$ while explicitly partisan
redistricting produces $|EG| = 0.08$ constitutes quantitative evidence
that compact, population-balanced districts are inherently fairer than
human-drawn ones---not because the algorithm tries to be fair, but
because it cannot try to be unfair.
